% Configurazione
\documentclass{article}

\usepackage{titling} % Required for inserting the subtitle
\usepackage{graphicx} % Required for inserting images
\usepackage{tabularx} % Per l'ambiente tabularx (tabelle)
\usepackage{calc} % Sempre per le tabelle
\usepackage[hidelinks]{hyperref} % Per i collegamenti ipertestuali, ad esempio sulla table of contents
\usepackage[italian]{babel} % Per la lingua italiana nelle scritte automatiche
\usepackage{xcolor} % Per colorare il testo
\usepackage{colortbl} % Per colorare le celle delle tabelle
\usepackage{lipsum} % Per generare lorem ipsum
\usepackage[normalem]{ulem} % Per sottolineare il testo
\usepackage{array} % Per la visualizzazione fluttuante di array di domande e risposte
\usepackage{ragged2e} % Pacchetto necessario per \justifying che giustifica il testo di tabelle

\newcommand{\ulhref}[2]{\href{#1}{\uline{#2}}} % Nuovo comando per sottolineare i link
\newcommand{\ulref}[1]{\uline{\ref{#1}}} % Nuovo comando per sottolineare i collegamenti a immagini
\setlength{\parindent}{0pt} % Rimuove il rientro automatico dei paragrafi

\graphicspath{ {immagini/} {../../shared/immagini/} }


% Struttura
\begin{document}

\begin{minipage}{0.4\textwidth}
    \includegraphics[width=0.6\textwidth]{Logo Universita di Padova.png}
\end{minipage}
\begin{minipage}{0.55\textwidth}
    \textcolor{red}{\textbf{Università degli Studi di Padova}} \\
    \textcolor{red}{Laurea: Informatica} \\
    \textcolor{red}{Corso: Ingegneria del Software} \\
    \textcolor{red}{Anno Accademico: 2024/2025}
\end{minipage}

\begin{minipage}{0.4\textwidth}
    \includegraphics[width=0.6\textwidth]{Logo SWEg.png}
\end{minipage}
\begin{minipage}{0.55\textwidth}
    \textbf{Gruppo: SWEg Labs} \\
    Email: \textsf{gruppo.sweg@gmail.com}
\end{minipage}

\vspace{2cm}

{
\centering
\Huge\bfseries Lettera di Presentazione\par
}

\vspace{2cm}
\section{Informazioni generali}

\begin{itemize}
    \item \textbf{Tipo di riunione:} esterna con proponenti
    \item \textbf{Luogo:} meeting zoom
    \item \textbf{Data:} 22/10/2024
    \item \textbf{Ora inizio:} 17:00
    \item \textbf{Ora fine:} 17:40
    \item \textbf{Responsabile:} Stefani Riccardo
    \item \textbf{Notaio:} Verzotto Davide
    \item \textbf{Partecipanti:}
    \begin{itemize}
        \renewcommand{\labelitemii}{--}
        \item Bolognini Federica
        \item Fantinato Michael
        \item Loat Giacomo
        \item Righetto Filippo
        \item Stefani Riccardo
        \item Verzotto Davide
        \item Carlesso Gianluca (\emph{Ergon Informatica Srl})
        \item Tieppo Anna (\emph{Ergon Informatica Srl})
    \end{itemize}
\end{itemize}

\newpage
\section*{Registro delle modifiche}

\begin{table}[h]
    \centering
    \begin{tabular}{|l|l|l|p{5cm}|}
        \hline
        \rowcolor[gray]{0.9}
        \textbf{Vers.} & \textbf{Data} & \textbf{Autore} & \textbf{Descrizione}\\
        \hline
        1.0 & ... & ... & Approvazione del documento\\
        \hline
        0.9 & ... & ... & Verifica del documento\\
        \hline
        0.7 & ... & ... & ...\\
        \hline
        0.6 & ... & ... & ...\\
        \hline
        0.5 & ... & ... & ...\\
        \hline
        0.2 & 25-10-2024 & Riccardo Stefani & Scrittura della motivazione della scelta e della conclusione\\
        \hline
        0.1 & 25-10-2024 & Riccardo Stefani & Creazione del documento\\
        \hline
    \end{tabular}
\end{table}

\newpage

\section{Introduzione}

In questo documento viene presentata l'analisi dei capitolati d'appalto che abbiamo svolto al fine di valutarne la fattibilità e scegliere il più adatto 
per il gruppo di lavoro \emph{SWEg Labs}. Verranno esposte in maniera chiara le motivazioni alla base della scelta di candidarci per il capitolato 
\emph{BuddyBot}, e altrettanto verranno analizzate le altre due proposte che avevamo fortemente considerato, \emph{LLM: Assistente digitale} e 
\emph{Vimar GENIALE}. Verranno altresì elencati i pro e i contro di ciascuno dei capitolati rimanenti per argomentare la decisione di averli scartati.
\newpage
\section{Resoconto degli incontri}

Example
\section{Motivazione della scelta}

La scelta finale è ricaduta su questo capitolato per diversi motivi:
\begin{itemize}
    \item \textit{\textbf{Disponibilità dell'azienda}}: L'azienda si è dimostrata da subito molto disponibile; le risposte alle mail sono arrivate in brevissimo tempo e si è subito manifestata la volontà di organizzare un incontro. Inoltre, si sono proposti di fornire anche supporto tecnico durante lo sviluppo, qualora ce ne fosse bisogno.
    \item \textit{\textbf{Interesse}}: Il prodotto finale ha suscitato grande interesse nel gruppo, soprattutto per le tecnologie con cui prevediamo di lavorare per realizzarlo. In particolare, ci è sembrato il capitolato con il miglior utilizzo dell'IA.
    \item \textit{\textbf{Fattibilità del progetto}}: Confrontando le richieste tra i vari capitolati, BuddyBot ci è sembrato il miglior compromesso tra difficoltà e interesse.
    \item \textit{\textbf{Libertà tecnologica}}: La libertà lasciata nella scelta delle tecnologie ci è sembrata molto vantaggiosa, permettendoci di utilizzare tecnologie con le quali alcuni membri del gruppo hanno già esperienza.
\end{itemize}

Visti i motivi sopra citati e non trovando in nessun altro capitolato una combinazione equivalente, 
il gruppo ha deciso di scegliere BuddyBot come progetto per cui candidarsi.
\newpage
Come riportato nel documento \emph{Preventivo dei costi ed impegni orari}, il costo stimato del prodotto ammonta a \textbf{**.000€};
si prevede inoltre di consegnarlo entro il \emph{***}.

\vspace{0.5cm}

Di seguito i componenti del gruppo:

\begin{table}[h]
    \centering
    \begin{tabular}{|c|c|c|}
        \hline
        \rowcolor[gray]{0.9}
        \textbf{Cognome} & \textbf{Nome} & \textbf{Matricola}\\
        \hline
        Bolognini & Federica & 2011881 \\
        \hline
        Fantinato & Michael & 2043672 \\
        \hline
        Loat & Giacomo & 2077677 \\
        \hline
        Righetto & Filippo & 2046428 \\
        \hline
        Stefani & Riccardo & 2068225 \\
        \hline
        Verzotto & Davide & 2043679 \\
        \hline
    \end{tabular}
\end{table}

\vspace{0.5cm}

Cordiali saluti,
\begin{flushright}
    Il gruppo \emph{SWEg Labs}
\end{flushright}


\end{document}
